\chapter*{Conclusiones y Próximos pasos}

Este manual de matemáticas para ICPC buscó que los alumnos pudieran acercarse a temas matemáticos a través de retos concretos. Como parte de la validación de los problemas de autoría, se compartió el grupo de Codeforces con un grupo reducido de compañeros con experiencia en ICPC. El problema \textit{Proyecto Hail Mary} fue el más resuelto (4 de 5 participantes), aunque con un número considerable de intentos fallidos. Si bien la muestra es pequeña, el patrón general indica que el problema requiere un pequeño análisis antes de poder aplicar la herramienta matemática subyacente directamente. Fue gratificante saber que la retroalimentación de los problemas de autoría propia fue positiva, varios alumnos reportaron
haberse mantenido entretenidos y determinados a resolverlos, incluso cuando implicó aprender técnicas que no conocían.

Una de las motivaciones detrás de este manual fue mostrar que el contenido de la licenciatura en matemáticas no vive aislado de la programación competitiva; al contrario, gran parte de las herramientas que aquí se desarrollan parten de temas que cualquier estudiante de la Facultad ya ve en su formación. Conocimientos de las materias de tronco común de la facultad fundamentan una gran parte del texto, un poco de teoría de grupos aparece al tocar el lema de Burnside, incluso Variable Compleja se asoma en la intuición detrás de la FFT, donde las raíces de la unidad complejas son el punto de partida antes de trasladar la idea al mundo modular con NTT. Y por supuesto, Teoría de Números y Teoría de Gráficas son, casi en su totalidad, prerrequisitos directos de la mitad de los temas cubiertos aquí. Ojalá este manual ayude a construir un puente, que un problema de Codeforces sea una oportunidad para repasar un teorema visto en clase, y que un teorema visto en clase sea una herramienta más en la mochila al momento de leer
un problema nuevo.

Escribir este manual fue en sí mismo un proceso de aprendizaje. Cada sección obligó a repensar ideas que parecían claras, a buscar la explicación más honesta posible y a encontrar los ejemplos que mejor ilustraran cada concepto. También confirmó lo abrumador que puede resultar la cantidad de recursos disponibles en línea: blogs, editoriales,tutoriales, libros, videos; dispersos, en distintos idiomas, con distintos niveles de rigor o diferentes enfoques. Este manual intentó ser un granito de arena en ese panorama: un recurso en español, con enfoque matemático, pensado para el contexto de la Facultad de Ciencias y las competencias nacionales y regionales del ICPC. Si le es útil a aunque sea un alumno, habrá valido la pena.

\section*{Trabajo futuro}

Hay varias direcciones naturales para extender este trabajo:

\textbf{Nuevas áreas temáticas.} Quedan fuera del alcance de esta
edición áreas importantes como probabilidad y teoría de juegos, las cuales son muy usadas en ICPC. A su vez, en áreas ya incluidas en este manual se podrían agregar temas más avanzados: flujo de costo mínimo para la sección de gráficas, juegos combinatorios como Hackenbush y Nim para combinatoria, triangulaciones de Delaunay y diagramas de Voronoi para geometría (con aplicaciones directas en problemas de geometría computacional más complejos), y pruebas de primalidad como Miller-Rabin para teoría de números, entre otros.

\textbf{Sección de problemas mixtos.} Una limitación del formato
actual es que los problemas están organizados por sección, lo que
puede llevar al lector a esperar de antemano qué herramienta usar. Se propone agregar una sección final con problemas sin clasificar, donde la dificultad radique precisamente en identificar qué técnicas combinar, más cercano a la experiencia real de una competencia.

\textbf{Más problemas de autoría propia.} Los problemas marcados con $\filledstar$ fueron muy bien recibidos por los lectores por intentar incluir narraciones cercanas a ellos que los motivara a resolverlos.
Extender la lista, especialmente con problemas que combinen varias áreas del manual, sería una contribución valiosa para una edición futura.

